\chapter{Topology of Two-Band Hamiltonians}
\chaptermark{Topology of Two-Band Hamiltonians}
\label{chap:topology-two-band-hamiltonians}

The previous chapters developed the algebraic machinery needed to reduce many
one-dimensional spin models to quadratic fermions and BdG momentum-space
blocks.  Once the Hamiltonian has been written as a family of finite-dimensional
matrices depending on momentum, exact solvability is not only a question of
finding the spectrum.  It also becomes meaningful to ask whether the family of
Bloch or BdG Hamiltonians can be continuously deformed into a trivial one while
keeping the bulk gap open and preserving the required symmetries.

This chapter gives the minimal topological toolkit needed for the later chapters
on the Kitaev chain, the SSH chain, and the XY/TFIM mother model.  The central
object is a two-band Hamiltonian
\begin{equation}
\label{eq:topo-two-band-general}
    H(k)=d_0(k)\id+\bm d(k)\cdot\bm\sigma,
    \qquad
    \bm d(k)=(d_x(k),d_y(k),d_z(k)),
\end{equation}
where $k$ belongs to the Brillouin zone.  For BdG Hamiltonians the Pauli
matrices act in Nambu space, while for ordinary band Hamiltonians they may act
in sublattice, orbital, or pseudospin space.  The same formulas therefore apply
to two physically different but mathematically parallel settings:
\begin{equation}
\label{eq:topo-two-realizations}
\begin{array}{ccl}
\text{SSH chain} &:& \text{sublattice two-band Hamiltonian},\\[2pt]
\text{Kitaev chain} &:& \text{BdG two-band Hamiltonian in Nambu space}.
\end{array}
\end{equation}
The distinction between these two realizations is important: the Pauli matrices
have the same algebra, but the symmetries and the meaning of the two components
of the spinor are different.

\section{Two-band Hamiltonians and maps to the Bloch sphere}
\label{sec:topo-two-band-bloch-sphere}

The spectrum of \cref{eq:topo-two-band-general} is
\begin{equation}
\label{eq:topo-two-band-spectrum}
    E_\pm(k)=d_0(k)\pm \abs{\bm d(k)}.
\end{equation}
Thus a two-band model is gapped at momentum $k$ precisely when
\begin{equation}
\label{eq:topo-gap-condition}
    \abs{\bm d(k)}\neq 0.
\end{equation}
For topology the scalar term $d_0(k)\id$ is usually irrelevant: it shifts the two
energies by the same amount but does not change the eigenvectors.  If the model
is gapped over the entire Brillouin zone, one may define the normalized vector
\begin{equation}
\label{eq:topo-normalized-d-vector}
    \widehat{\bm d}(k)=\frac{\bm d(k)}{\abs{\bm d(k)}}\in S^2 .
\end{equation}
Therefore a gapped two-band Hamiltonian defines a map
\begin{equation}
\label{eq:topo-map-bz-to-sphere}
    \widehat{\bm d}:\mathrm{BZ}\longrightarrow S^2.
\end{equation}
A topological phase is a homotopy class of such maps, possibly restricted by
symmetry.

The projectors onto the upper and lower bands of $\bm d\cdot\bm\sigma$ are
\begin{equation}
\label{eq:topo-band-projectors}
    P_+(k)=\frac{1+\widehat{\bm d}(k)\cdot\bm\sigma}{2},
    \qquad
    P_-(k)=\frac{1-\widehat{\bm d}(k)\cdot\bm\sigma}{2}.
\end{equation}
If the lower band is occupied, the occupied-band geometry is completely encoded
in $P_-(k)$.

\begin{principle}[Gap closing is the only way to change topology]
For a two-band Hamiltonian, the topological invariant can change only if there
exists a momentum $k_0$ such that
\begin{equation}
    \bm d(k_0)=0.
\end{equation}
At this point the two bands touch, the normalized vector
$\widehat{\bm d}(k)$ is undefined, and the homotopy class of the map can change.
\end{principle}

This principle is the topological version of the gap-closing conditions already
encountered in the BdG chapter.  For example, the Kitaev-chain spectrum
\begin{equation}
    E(k)=\sqrt{\xi(k)^2+\abs{\Delta(k)}^2}
\end{equation}
can close only when both $\xi(k)$ and $\Delta(k)$ vanish.  Since the nearest-neighbor
$p$-wave gap is proportional to $\sin k$, the possible gap closings occur at
$k=0,\pi$.

\section{Berry connection, Berry curvature, and Chern number}
\label{sec:topo-berry-curvature-chern}

Although much of this chapter focuses on one-dimensional models, the general
two-band language is most transparent if we first state the two-dimensional
Berry-curvature formula.  Let $\ket{u_-(\bm k)}$ be a normalized eigenvector of
the lower band.  The Berry connection and Berry curvature are
\begin{equation}
\label{eq:topo-berry-connection-curvature}
    A_a(\bm k)=\ii\braket{u_-(\bm k)}{\partial_{k_a}u_-(\bm k)},
    \qquad
    F_{xy}=\partial_{k_x}A_y-\partial_{k_y}A_x .
\end{equation}
Equivalently, in a manifestly gauge-invariant form,
\begin{equation}
\label{eq:topo-projector-curvature}
    F_{xy}(\bm k)
    =
    \ii\Tr\Bigl(P_-(\bm k)\bigl[\partial_{k_x}P_-(\bm k),
    \partial_{k_y}P_-(\bm k)\bigr]\Bigr).
\end{equation}
Substituting \cref{eq:topo-band-projectors} gives the standard two-band formula
\begin{equation}
\label{eq:topo-two-band-berry-curvature}
    F_{xy}(\bm k)
    =
    \frac12\widehat{\bm d}\cdot
    \left(
        \partial_{k_x}\widehat{\bm d}
        \times
        \partial_{k_y}\widehat{\bm d}
    \right).
\end{equation}
The first Chern number of the occupied band is
\begin{equation}
\label{eq:topo-chern-number}
    C
    =
    \frac{1}{2\pi}\int_{\mathrm{BZ}} F_{xy}(\bm k)\,\dd^2k
    =
    \frac{1}{4\pi}\int_{\mathrm{BZ}}
    \widehat{\bm d}\cdot
    \left(
        \partial_{k_x}\widehat{\bm d}
        \times
        \partial_{k_y}\widehat{\bm d}
    \right)\dd^2k .
\end{equation}
Geometrically, $C$ counts how many times the Brillouin-zone torus wraps the unit
sphere under the map $\widehat{\bm d}$.

\begin{example}[A standard lattice Chern-insulator block]
\label{ex:topo-qwz-model}
Consider
\begin{equation}
\label{eq:topo-qwz-d-vector}
    \bm d(\bm k)=
    \bigl(
        \sin k_x,
        \sin k_y,
        m+\cos k_x+\cos k_y
    \bigr).
\end{equation}
The gap can close only at the four high-symmetry points
$(0,0)$, $(\pi,0)$, $(0,\pi)$, and $(\pi,\pi)$, giving critical values
$m=-2,0,2$.  With the orientation convention in
\cref{eq:topo-chern-number}, the occupied-band Chern number is
\begin{equation}
\label{eq:topo-qwz-chern-values}
    C=\begin{cases}
    0, & m>2,\\
    -1, & 0<m<2,\\
    +1, & -2<m<0,\\
    0, & m<-2.
    \end{cases}
\end{equation}
This example is not one of the main solvable chains of these notes, but it is a
useful reference point: in two dimensions the full sphere $S^2$ can be wrapped,
whereas in one-dimensional chiral models topology usually comes from restricting
$\widehat{\bm d}$ to an equator.
\end{example}

\section{Why one-dimensional two-band topology needs symmetry}
\label{sec:topo-1d-needs-symmetry}

For a generic one-dimensional two-band Hamiltonian, the map is
\begin{equation}
    S^1_{\rm BZ}\longrightarrow S^2 .
\end{equation}
Since $\pi_1(S^2)=0$, every loop on the sphere can be contracted to a point.
Therefore a generic gapped one-dimensional two-band Hamiltonian has no stable
integer winding number.

A nontrivial one-dimensional invariant appears when a symmetry forces
$\widehat{\bm d}(k)$ to lie on a great circle of the sphere.  The most important
case for these notes is chiral symmetry.

\begin{definition}[Chiral symmetry]
A two-band Hamiltonian has chiral symmetry if there exists a unitary matrix
$S$ with $S^2=\id$ such that
\begin{equation}
\label{eq:topo-chiral-symmetry-definition}
    S H(k) S^{-1}=-H(k).
\end{equation}
For the conventional choice $S=\sigma_z$, chiral symmetry forbids the identity
and $\sigma_z$ terms, so that
\begin{equation}
\label{eq:topo-chiral-two-band}
    H(k)=d_x(k)\sigma_x+d_y(k)\sigma_y.
\end{equation}
\end{definition}

The vector $\bm d(k)$ now lies in the $d_x$--$d_y$ plane.  If the system is
gapped, the complex number
\begin{equation}
\label{eq:topo-q-definition}
    q(k)=d_x(k)+\ii d_y(k)
\end{equation}
never vanishes.  Thus $q(k)$ defines a map
\begin{equation}
\label{eq:topo-q-map}
    S^1_{\rm BZ}\longrightarrow \CC\setminus\{0\}\simeq S^1 .
\end{equation}
This map is classified by an integer winding number.

\section{Winding number of chiral two-band Hamiltonians}
\label{sec:topo-winding-number}

For the chiral Hamiltonian
\begin{equation}
\label{eq:topo-chiral-hamiltonian}
    H(k)=d_x(k)\sigma_x+d_y(k)\sigma_y
    =
    \begin{pmatrix}
        0 & q(k)^*\\
        q(k) & 0
    \end{pmatrix},
    \qquad
    q(k)=d_x(k)+\ii d_y(k),
\end{equation}
where the second equality follows from the standard Pauli matrices.  The winding
number is
\begin{equation}
\label{eq:topo-winding-number}
    \nu
    =
    \frac{1}{2\pi\ii}\int_{\rm BZ}\dd k\,\partial_k\log q(k).
\end{equation}
Equivalently, writing
\begin{equation}
    q(k)=\abs{q(k)}\ee^{\ii\varphi(k)},
\end{equation}
one obtains
\begin{equation}
\label{eq:topo-winding-angle}
    \nu
    =
    \frac{1}{2\pi}\int_{\rm BZ}\dd k\,\partial_k\varphi(k)
    =
    \frac{\varphi(2\pi)-\varphi(0)}{2\pi}\in\ZZ .
\end{equation}
Thus $\nu$ counts how many times the curve
\begin{equation}
    k\longmapsto (d_x(k),d_y(k))
\end{equation}
winds around the origin.

\begin{proposition}[Topological stability of the winding number]
If $q_s(k)$ is a smooth one-parameter family with $q_s(k)\neq 0$ for all
$k$ and all $s\in[0,1]$, then the winding number \cref{eq:topo-winding-number}
is independent of $s$.
\end{proposition}

\begin{proof}
The $s$-derivative of the winding number is
\begin{equation}
    \partial_s\nu
    =
    \frac{1}{2\pi\ii}\int_{\rm BZ}\dd k\,
    \partial_k\left(q^{-1}\partial_s q\right).
\end{equation}
Since the Brillouin zone is a circle and $q^{-1}\partial_s q$ is single-valued
when $q$ never vanishes, the integral of the total derivative is zero.
\end{proof}

The winding number can therefore change only when $q(k_0)=0$ for some momentum.
This is the one-dimensional chiral version of the general gap-closing principle.

\section{Zak phase and its relation to winding}
\label{sec:topo-zak-phase}

For a one-dimensional occupied band, the Berry phase over the Brillouin zone is
called the Zak phase:
\begin{equation}
\label{eq:topo-zak-phase-definition}
    \gamma_{\rm Zak}
    =
    \oint_{\rm BZ}\dd k\,
    \ii\braket{u_-(k)}{\partial_k u_-(k)}
    \quad\mod 2\pi .
\end{equation}
For a chiral two-band Hamiltonian, one may choose the lower-band eigenvector as
\begin{equation}
\label{eq:topo-chiral-eigenvector}
    \ket{u_-(k)}
    =
    \frac{1}{\sqrt{2}}
    \begin{pmatrix}
        -\ee^{-\ii\varphi(k)}\\
        1
    \end{pmatrix},
    \qquad
    q(k)=\abs{q(k)}\ee^{\ii\varphi(k)}.
\end{equation}
This gives
\begin{equation}
\label{eq:topo-zak-winding-relation}
    \gamma_{\rm Zak}
    =
    +\frac12\int_{\rm BZ}\dd k\,\partial_k\varphi(k)
    =
    +\pi\nu
    \quad\mod 2\pi .
\end{equation}
The sign depends on the eigenvector and Pauli-matrix convention, but the
quantized statement is invariant:
\begin{equation}
\label{eq:topo-zak-quantized-statement}
    \gamma_{\rm Zak}=0\ \text{or}\ \pi\quad\mod 2\pi
\end{equation}
for a chiral or inversion-symmetric two-band chain.  Thus the parity of the
winding number is visible as a Berry phase.

\section{SSH chain as a chiral two-band model}
\label{sec:topo-ssh-chain}

The Su--Schrieffer--Heeger chain is the canonical ordinary-band example.  It has
two sublattices $A$ and $B$ in each unit cell, with intracell hopping $t_1$ and
intercell hopping $t_2$:
\begin{equation}
\label{eq:topo-ssh-real-space}
    H_{\rm SSH}
    =
    \sum_j
    \left(
        t_1 a_j^\dagger b_j
        +t_2 a_{j+1}^\dagger b_j
        +\text{h.c.}
    \right).
\end{equation}
Using the Bloch spinor
\begin{equation}
    \Psi_k=\begin{pmatrix}a_k\\ b_k\end{pmatrix},
\end{equation}
the Hamiltonian becomes
\begin{equation}
\label{eq:topo-ssh-bloch-hamiltonian}
    H_{\rm SSH}
    =
    \sum_k\Psi_k^\dagger
    \begin{pmatrix}
        0 & q_{\rm SSH}(k)^*\\
        q_{\rm SSH}(k) & 0
    \end{pmatrix}
    \Psi_k,
    \qquad
    q_{\rm SSH}(k)=t_1+t_2\ee^{\ii k}.
\end{equation}
Equivalently,
\begin{equation}
\label{eq:topo-ssh-d-vector}
    d_x(k)=t_1+t_2\cos k,
    \qquad
    d_y(k)=t_2\sin k,
    \qquad
    d_z(k)=0.
\end{equation}
The curve $q_{\rm SSH}(k)$ is a circle of radius $\abs{t_2}$ centered at $t_1$ on
the real axis.  It winds around the origin if and only if the origin lies inside
this circle.  For $t_1,t_2>0$ and the convention in
\cref{eq:topo-ssh-bloch-hamiltonian},
\begin{equation}
\label{eq:topo-ssh-winding-result}
    \nu_{\rm SSH}
    =
    \begin{cases}
    0, & \abs{t_1}>\abs{t_2},\\
    1, & \abs{t_2}>\abs{t_1}.
    \end{cases}
\end{equation}
The transition occurs at $\abs{t_1}=\abs{t_2}$, where the gap closes at
$k=\pi$ for $t_1=t_2$.

\begin{remark}[Unit-cell convention and the sign of winding]
Changing the unit-cell convention multiplies $q(k)$ by a phase
$\ee^{\ii n k}$.  This shifts the integer assigned to a particular Bloch
basis.  The robust physical statement for the SSH chain with a fixed boundary
termination is whether the dimerization leaves unpaired edge sites.  In the
convention above, $\abs{t_2}>\abs{t_1}$ is the topological phase for the usual
termination with weak intracell bonds at the boundary.
\end{remark}

\subsection{Edge zero modes}
\label{subsec:topo-ssh-edge-zero-modes}

For an open chain in the topological regime $\abs{t_2}>\abs{t_1}$, the left-edge
zero mode has support only on one sublattice.  A zero-energy ansatz of the form
\begin{equation}
\label{eq:topo-ssh-left-edge-ansatz}
    \ket{\psi_L}=\sum_{j\geq 1}\alpha_j a_j^\dagger\ket{0}
\end{equation}
obeys the recursion
\begin{equation}
\label{eq:topo-ssh-edge-recursion}
    t_1\alpha_j+t_2\alpha_{j+1}=0,
    \qquad
    \alpha_{j+1}=-\frac{t_1}{t_2}\alpha_j .
\end{equation}
Therefore
\begin{equation}
\label{eq:topo-ssh-edge-wavefunction}
    \alpha_j=\left(-\frac{t_1}{t_2}\right)^{j-1}\alpha_1,
\end{equation}
which is normalizable precisely when $\abs{t_1/t_2}<1$.  The bulk winding and
the boundary zero mode therefore encode the same phase distinction.

\section{BdG topology: Kitaev chain and Majorana zero modes}
\label{sec:topo-bdg-kitaev}

The Kitaev chain is also a two-band Hamiltonian, but the two components of the
spinor are not sublattices.  They are particle and hole components of a Nambu
spinor.  In the convention of the BdG chapter,
\begin{equation}
\label{eq:topo-kitaev-bdg-block}
    \mathcal{H}_{\rm K}(k)
    =
    \begin{pmatrix}
        \xi(k) & \Delta_\phi(k)\\
        \Delta_\phi(k)^* & -\xi(k)
    \end{pmatrix},
    \qquad
    \xi(k)=-\mu-2t\cos k,
\end{equation}
with
\begin{equation}
\label{eq:topo-kitaev-gap-function}
    \Delta_\phi(k)
    =
    2\ii\abs{\Delta}\ee^{\ii(\theta_\Delta-2\phi)}\sin k .
\end{equation}
The uniform phase $\theta_\Delta$ and the Fourier/Nambu gauge phase $\phi$ enter
only through the combination $\theta_\Delta-2\phi$.  For a single uniform chain,
one may choose
\begin{equation}
    \phi=\frac{\theta_\Delta}{2}
\end{equation}
so that
\begin{equation}
\label{eq:topo-kitaev-real-gap-gauge}
    \Delta_\phi(k)=2\ii\abs{\Delta}\sin k.
\end{equation}
In this real pairing gauge,
\begin{equation}
\label{eq:topo-kitaev-d-vector}
    \mathcal{H}_{\rm K}(k)
    =
    d_y(k)\tau_y+d_z(k)\tau_z,
    \qquad
    d_y(k)=-2\abs{\Delta}\sin k,
    \qquad
    d_z(k)=\xi(k).
\end{equation}
Here the Pauli matrices $\tau_a$ act in Nambu space.  Since the vector lies in a
plane, the real Kitaev chain has a winding-number description analogous to a
chiral two-band model.

A convenient complex function is
\begin{equation}
\label{eq:topo-kitaev-q-function}
    q_{\rm K}(k)=\xi(k)+2\ii\abs{\Delta}\sin k
    =-\mu-2t\cos k+2\ii\abs{\Delta}\sin k.
\end{equation}
Its winding number is
\begin{equation}
\label{eq:topo-kitaev-bdi-winding}
    \nu_{\rm BDI}
    =
    \frac{1}{2\pi\ii}\int_{\rm BZ}\dd k\,
    \partial_k\log q_{\rm K}(k).
\end{equation}
Geometrically, $q_{\rm K}(k)$ traces an ellipse centered at $-\mu$ on the real
axis.  The origin lies inside this ellipse if and only if
\begin{equation}
\label{eq:topo-kitaev-topological-condition}
    \abs{\mu}<2\abs{t}.
\end{equation}
Thus the Kitaev chain is topological for $\abs{\mu}<2\abs{t}$ and trivial for
$\abs{\mu}>2\abs{t}$, assuming $\Delta\neq0$.

\subsection{Class D Pfaffian invariant}
\label{subsec:topo-class-d-pfaffian}

The real gauge above makes an additional chiral symmetry manifest and places the
nearest-neighbor Kitaev chain in class BDI.  More generally, a superconducting
BdG Hamiltonian always has particle-hole symmetry.  In the Nambu convention
used earlier,
\begin{equation}
\label{eq:topo-phs-condition}
    \tau_x\mathcal{H}_{\rm BdG}(k)^*\tau_x=-\mathcal{H}_{\rm BdG}(-k).
\end{equation}
Without a real gauge or time-reversal-like symmetry, the one-dimensional
superconductor is in class D and has a $\ZZ_2$ invariant rather than an integer
winding number.

For a two-band Kitaev chain the class-D invariant reduces to the sign of the
normal-state term at the particle-hole-invariant momenta $k=0,\pi$, because the
odd-parity pairing vanishes there:
\begin{equation}
\label{eq:topo-kitaev-class-d-invariant}
    \mathcal{M}
    =
    \sgn\bigl[\xi(0)\xi(\pi)\bigr]
    =
    \sgn\bigl[(-\mu-2t)(-\mu+2t)\bigr]
    =
    \sgn(\mu^2-4t^2).
\end{equation}
The nontrivial phase is
\begin{equation}
\label{eq:topo-class-d-nontrivial-condition}
    \mathcal{M}=-1
    \quad\Longleftrightarrow\quad
    \abs{\mu}<2\abs{t}.
\end{equation}
This is the $\ZZ_2$ remnant of the integer BDI winding.  Uniform complex phases
of $\Delta$ do not change \cref{eq:topo-kitaev-class-d-invariant}; they can be
removed by a global Nambu gauge choice.  Relative superconducting phases,
spatially varying phases, or fluxes that cannot be removed globally can reduce
the manifest symmetry to class D, but the $\ZZ_2$ invariant remains well-defined
as long as particle-hole symmetry and the bulk gap persist.

\subsection{Exactly solvable Majorana limit}
\label{subsec:topo-majorana-edge-limit}

The boundary meaning of \cref{eq:topo-kitaev-topological-condition} is clearest
at the exactly solvable point
\begin{equation}
\label{eq:topo-kitaev-sweet-spot}
    \mu=0,
    \qquad
    t=\abs{\Delta}>0,
\end{equation}
with open boundary conditions.  Using Majorana operators
\begin{equation}
\label{eq:topo-majorana-definitions}
    a_{2j-1}=c_j+c_j^\dagger,
    \qquad
    a_{2j}=\ii(c_j-c_j^\dagger),
\end{equation}
the Hamiltonian becomes, up to an overall sign convention fixed by the phase of
$\Delta$,
\begin{equation}
\label{eq:topo-kitaev-majorana-sweet-spot}
    H_{\rm K}
    =
    -\ii t\sum_{j=1}^{L-1}a_{2j}a_{2j+1}.
\end{equation}
All bulk Majoranas are paired on bonds, while $a_1$ and $a_{2L}$ are absent from
the Hamiltonian.  These two unpaired Majorana operators form a nonlocal boundary
fermion
\begin{equation}
\label{eq:topo-edge-fermion}
    f_{\rm edge}=\frac{a_1+\ii a_{2L}}{2}.
\end{equation}
The two occupations of $f_{\rm edge}$ are degenerate in the thermodynamic limit.
Away from the exactly solvable point but within the same gapped topological
phase, the Majorana zero modes remain exponentially localized at the edges.

\begin{principle}[Bulk-boundary correspondence in one-dimensional two-band models]
A nontrivial bulk winding or class-D $\ZZ_2$ invariant predicts protected
boundary zero modes, provided the boundary termination respects the protecting
symmetry and the bulk gap remains open.
\end{principle}

For the SSH chain the boundary zero modes are ordinary complex fermion states
localized on sublattices.  For the Kitaev chain they are Majorana zero modes in
Nambu space.  The algebraic similarity of the two-band Hamiltonians should not
obscure this physical difference.

\section{Relation to the TFIM and the Majorana mass sign}
\label{sec:topo-relation-to-tfim}

The duality chapter showed that the transverse-field Ising chain separates two
phases exchanged by Kramers--Wannier duality.  The Jordan--Wigner and BdG
chapters showed that the same model is equivalent to a quadratic Majorana chain.
Topology supplies a third language for the same transition.

Near a gap-closing momentum, a one-dimensional BdG Hamiltonian takes the
continuum Dirac/Majorana form
\begin{equation}
\label{eq:topo-continuum-majorana-dirac}
    \mathcal{H}(q)=v q\,\tau_y+m\,\tau_z+\text{higher-order terms},
\end{equation}
where $q$ measures momentum relative to the gap-closing point.  A topological
transition occurs when the mass $m$ changes sign.  For the Kitaev chain,
\begin{equation}
\label{eq:topo-kitaev-masses}
    m_0=\xi(0)=-\mu-2t,
    \qquad
    m_\pi=\xi(\pi)=-\mu+2t.
\end{equation}
The class-D invariant is simply the relative sign of these two masses:
\begin{equation}
\label{eq:topo-mass-sign-invariant}
    \mathcal{M}=\sgn(m_0m_\pi).
\end{equation}
The nontrivial phase has opposite signs at $k=0$ and $k=\pi$.

In the TFIM image, the topological Majorana phase corresponds to the ordered
Ising phase for open chains.  The edge Majoranas encode the same twofold
structure that appears as symmetry breaking in the spin variables.  Thus the
following descriptions are equivalent ways of describing the same physics:
\begin{equation}
\label{eq:topo-tfim-equivalent-descriptions}
\begin{array}{ccl}
\text{spin language} &:& \text{ordered versus disordered Ising phase},\\[2pt]
\text{duality language} &:& \text{order variable versus disorder variable},\\[2pt]
\text{fermion language} &:& \text{trivial versus topological Majorana chain},\\[2pt]
\text{continuum language} &:& \text{sign of a Majorana mass}.
\end{array}
\end{equation}

\section{Symmetry classes relevant to the mother models}
\label{sec:topo-symmetry-classes}

The full tenfold classification is beyond the scope of these notes, but the
following restricted table is sufficient for the solvable models treated here.

\begin{center}
\small
\begin{tabular}{L{0.15\textwidth}L{0.22\textwidth}L{0.21\textwidth}L{0.28\textwidth}}
\toprule
Class & Symmetry content & 1D invariant & Representative model in these notes \\
\midrule
AIII & Chiral symmetry only & $\ZZ$ winding & SSH chain with chiral sublattice symmetry \\
BDI & Particle-hole plus real time-reversal-like symmetry, hence chiral symmetry & $\ZZ$ winding & Real Kitaev chain / JW image of XY-type chains in a real gauge \\
D & Particle-hole symmetry only & $\ZZ_2$ Pfaffian invariant & Kitaev chain with only BdG particle-hole symmetry manifest \\
A & No protecting symmetry & No stable 1D two-band invariant; $\ZZ$ Chern number in 2D & Two-dimensional Chern-insulator block \\
\bottomrule
\end{tabular}
\end{center}

The table should be read operationally.  The same matrix expression may belong
to different classes depending on what physical structure the Pauli matrices
represent and which perturbations are allowed.  For example, adding a sublattice
staggering term $m\sigma_z$ to the SSH chain breaks chiral symmetry and destroys
the integer winding classification.  Adding a perturbation to a BdG Hamiltonian
must still respect particle-hole redundancy, so the allowed deformations are
different.

\section{Gauge choices, phases, and what is physically invariant}
\label{sec:topo-gauge-phases}

The Fourier and BdG chapters emphasized that a constant phase in the Fourier
transform,
\begin{equation}
\label{eq:topo-fourier-gauge-reminder}
    c_k=\frac{\ee^{-\ii\phi}}{\sqrt{L}}\sum_j\ee^{-\ii k j}c_j,
\end{equation}
is a gauge convention.  In BdG language it rotates the Nambu spinor and changes
the phase of the gap function as
\begin{equation}
\label{eq:topo-gap-gauge-rotation}
    \Delta_\phi(k)=\ee^{-2\ii\phi}\Delta_0(k).
\end{equation}
Consequently the components $(d_x,d_y)$ of a BdG vector may rotate in the
transverse plane under a gauge change.  This rotation does not change the energy
spectrum, the existence of a gap, or the class-D invariant
\cref{eq:topo-kitaev-class-d-invariant}.

For an ordinary chiral band Hamiltonian such as the SSH chain, a change of Bloch
basis can also multiply $q(k)$ by a phase.  If the phase is periodic and smooth,
it does not change the physically meaningful distinction between phases with and
without protected boundary states for a fixed termination.  If the phase has the
form $\ee^{\ii n k}$, it reflects a shift of unit-cell convention and changes the
integer assigned to the winding by a convention-dependent offset.  Boundary
physics is recovered only after the unit cell and termination are specified.

The correct hierarchy is therefore
\begin{equation}
\label{eq:topo-invariant-hierarchy}
\begin{array}{ccl}
\text{basis-dependent} &:& d_x(k),d_y(k),d_z(k)\ \text{as components},\\[2pt]
\text{gauge-covariant} &:& q(k)\ \text{and its phase},\\[2pt]
\text{physically invariant} &: &
\begin{gathered}
\text{gap closings, Berry phases modulo }2\pi,\\
\text{winding differences, }\ZZ_2\text{ parity, and edge modes}
\end{gathered}.
\end{array}
\end{equation}
This distinction is especially important for the Kitaev chain with a complex
superconducting gap: a uniform complex phase is removable by a global gauge
choice, while a relative phase across a junction or a spatial phase twist can
represent physical superconducting phase bias or flux.

\section{Practical diagnostic formulas}
\label{sec:topo-practical-diagnostics}

For later use, we collect the most useful diagnostic formulas.

\subsection*{Chiral two-band chain}
For
\begin{equation}
    H(k)=d_x(k)\sigma_x+d_y(k)\sigma_y,
    \qquad
    q(k)=d_x(k)+\ii d_y(k),
\end{equation}
the integer invariant is
\begin{equation}
\label{eq:topo-summary-chiral-winding}
    \nu=\frac{1}{2\pi\ii}\oint_{\rm BZ}\dd k\,\partial_k\log q(k).
\end{equation}
A gap closing occurs when $q(k)=0$.

\subsection*{SSH chain}
For
\begin{equation}
    q_{\rm SSH}(k)=t_1+t_2\ee^{\ii k},
\end{equation}
the usual topological regime for the standard open-chain termination is
\begin{equation}
\label{eq:topo-summary-ssh}
    \abs{t_2}>\abs{t_1}.
\end{equation}

\subsection*{Kitaev chain}
For
\begin{equation}
    \xi(k)=-\mu-2t\cos k,
    \qquad
    \Delta(k)=2\ii\abs{\Delta}\sin k
\end{equation}
in a real pairing gauge, the BDI winding may be computed from
\begin{equation}
\label{eq:topo-summary-kitaev-winding}
    q_{\rm K}(k)=\xi(k)+2\ii\abs{\Delta}\sin k.
\end{equation}
The class-D invariant is
\begin{equation}
\label{eq:topo-summary-kitaev-z2}
    \mathcal{M}=\sgn\bigl[\xi(0)\xi(\pi)\bigr]
    =\sgn(\mu^2-4t^2).
\end{equation}
The topological Majorana phase is
\begin{equation}
\label{eq:topo-summary-kitaev-topological-phase}
    \mathcal{M}=-1
    \quad\Longleftrightarrow\quad
    \abs{\mu}<2\abs{t}.
\end{equation}

\subsection*{Two-dimensional two-band block}
For a two-dimensional occupied lower band,
\begin{equation}
\label{eq:topo-summary-chern}
    C
    =
    \frac{1}{4\pi}\int_{\rm BZ}
    \widehat{\bm d}\cdot
    \left(
        \partial_{k_x}\widehat{\bm d}
        \times
        \partial_{k_y}\widehat{\bm d}
    \right)\dd^2k.
\end{equation}

\section{Summary}
\label{sec:topo-summary}

The essential lesson of this chapter is that a solvable two-band Hamiltonian
contains more information than its dispersion.  The eigenvectors define a map
from momentum space to a target space, and the topology of this map constrains
possible phases and boundary states.

The main results are:
\begin{enumerate}[label=(\roman*)]
    \item A gapped two-band Hamiltonian defines a normalized vector
    $\widehat{\bm d}(k)$ and can change topology only when $\bm d(k)=0$.
    \item In one dimension, a generic two-band Hamiltonian is topologically
    trivial unless a symmetry restricts the target space.
    \item Chiral two-band chains have an integer winding number
    $\nu=(2\pi\ii)^{-1}\int\partial_k\log q(k)$.
    \item The SSH chain realizes this winding in sublattice space.
    \item The Kitaev chain realizes an analogous structure in Nambu space; in
    class D the robust invariant is the Pfaffian sign
    $\mathcal{M}=\sgn[\xi(0)\xi(\pi)]$.
    \item Majorana edge modes are the boundary manifestation of the nontrivial
    Kitaev-chain invariant.
    \item The TFIM transition, the Kitaev-chain topological transition, and the
    sign change of a continuum Majorana mass are different languages for the
    same solvable critical structure.
\end{enumerate}
